\section{Inference}

\todo{Substitutions are used to lift inference to propositional logics:
    \begin{itemize}
        \item Represent the formulas to be infered with the same computation network of atomic formulas.
        \item Represent the formulas with an entangled activation network (whereas before they were disentangled, since refering to different atomic sentences).
    \end{itemize}
}

We understand inference as manipulations in the activation network.
\begin{itemize}
    \item {\bf Unification}: Substitution of variables to represent multiple sentences by same activation network
    \item {\bf Propositional inference}: Manipulations in activation network, as for constraint networks
\end{itemize}

\subsection{Unification}

Unification of atomic sentences: Find a substitution to the variables, such that both sentences are represented by the same computation network.

Compressing of complex sentences:
\begin{itemize}
    \item Find a unification of the atomic sentences in both complex sentences
    \item Represent the substituted complex sentences by a constraction of the activation networks
\end{itemize}

\subsection{Modus ponens}

\begin{example}
    For example we have
    \begin{itemize}
        \item John knows everyone: Knows(John,z)
        \item People who know each other are friends: Knows(x,y) $\Rightarrow$ Friends(x,y)
    \end{itemize}
    The unification of Knows(John,z) and Knows(x,y) is given by the substitution $x\mapsto$ John and $y \mapsto z$.
    We apply the unification on the second rule and get Knows(John,y) $\Rightarrow$ Friends(John,y).


    Now we do modus ponens on the activation network and get Friends(John,z).
    \stantikzpic{
% From Knows(John,z)
        \drawsquaretensor{0}{4}{$\tbasis$}
        \draw (0,2) -- (0,3);

        % From Knows(x,y) $\Rightarrow$ Friends(x,y)
        \draw (3,3) rectangle (7,5);
        \drawtextnode{5}{4}{\corelabelsize $\Rightarrow$};
        \draw (0,2) to[bend right=20] (4,3);
        \draw (6,2) -- (6,3);

        \drawvariabledot{0}{2}
        \drawvariabledot{6}{2}

        %% SUBSTITUTION
        \draw [->-] (0,1) -- (0,2) node[midway,right] {\colorlabelsize $\indvariableof{Knows}$};
        \draw (-2,-1) rectangle (2,1);
        \drawtextnode{0}{0}{\corelabelsize $\bencodingof{Knows}$};

        \begin{scope}[shift={(6,0)}]
            \draw [->-] (0,1) -- (0,2) node[midway,right] {\colorlabelsize $\indvariableof{Friends}$};
            \draw (-2,-1) rectangle (2,1);
            \drawtextnode{0}{0}{\corelabelsize $\bencodingof{Friends}$};
        \end{scope}

        \draw[->-] (1.5,-4) to[bend right=-40] (-1.5,-1);
        \draw[->-] (1.5,-4) to[bend right=40] (4.5,-1);
        \drawvariabledot{1.5}{-4}
        \draw[->-] (1.5,-5) -- (1.5,-4) node[midway,right] {\colorlabelsize $\indvariableof{x}$};
        \drawsquaretensor{1.5}{-6}{$\onehotmapof{John}$}

        \draw[->-] (4.5,-4) to[bend right=-40] (1.5,-1);
        \draw[->-] (4.5,-4) to[bend right=40] (7.5,-1);
        \drawvariabledot{4.5}{-4}
        \draw[->-] (4.5,-5) -- (4.5,-4) node[midway,right] {\colorlabelsize $\indvariableof{z}$};
        \drawsquaretensor{4.5}{-6}{$\identity$}
        \draw[->-] (4.5,-8.5) -- (4.5,-7) node[midway,right] {\colorlabelsize $\indvariableof{z}$};

        \drawtextnode{13}{0}{${=}$}
    %%% RIGHT SIDE
        \begin{scope}[shift={(20,0)}]
            % From Knows(John,z)
            \drawsquaretensor{0}{4}{$\tbasis$}
            \draw (0,2) -- (0,3);

            % Modus ponens infered
            \drawsquaretensor{6}{4}{$\tbasis$}
            \draw (6,2) -- (6,3);

            \drawvariabledot{0}{2}
            \drawvariabledot{6}{2}

            %% SUBSTITUTION
            \draw [->-] (0,1) -- (0,2) node[midway,right] {\colorlabelsize $\indvariableof{Knows}$};
            \draw (-2,-1) rectangle (2,1);
            \drawtextnode{0}{0}{\corelabelsize $\bencodingof{Knows}$};

            \begin{scope}[shift={(6,0)}]
                \draw [->-] (0,1) -- (0,2) node[midway,right] {\colorlabelsize $\indvariableof{Friends}$};
                \draw (-2,-1) rectangle (2,1);
                \drawtextnode{0}{0}{\corelabelsize $\bencodingof{Friends}$};
            \end{scope}

            \draw[->-] (1.5,-4) to[bend right=-40] (-1.5,-1);
            \draw[->-] (1.5,-4) to[bend right=40] (4.5,-1);
            \drawvariabledot{1.5}{-4}
            \draw[->-] (1.5,-5) -- (1.5,-4) node[midway,right] {\colorlabelsize $\indvariableof{x}$};
            \drawsquaretensor{1.5}{-6}{$\onehotmapof{John}$}

            \draw[->-] (4.5,-4) to[bend right=-40] (1.5,-1);
            \draw[->-] (4.5,-4) to[bend right=40] (7.5,-1);
            \drawvariabledot{4.5}{-4}
            \draw[->-] (4.5,-5) -- (4.5,-4) node[midway,right] {\colorlabelsize $\indvariableof{z}$};
            \drawsquaretensor{4.5}{-6}{$\identity$}
            \draw[->-] (4.5,-8.5) -- (4.5,-7) node[midway,right] {\colorlabelsize $\indvariableof{z}$};
        \end{scope}
    }
\end{example}

\subsection{Contraction equations}

Quantified formulas understood are contraction equations.

Understanding them as hard constraints, one can infer parts of the substitution tensors, which holds for all supported worlds.

\begin{example}[Modus ponens on the substitution structure]
    The formula
    \begin{align*}
        \uniquantwrtof{\shortcatindices}{\formulaat{\indexedshortcatvariables}\impbincon\secformulaat{\indexedshortcatvariables}}
    \end{align*}
    corresponds with the contraction equation
    \stantikzpic{
        \drawtensorblock{2.5}{4}{3}{$\impbincon$}
        \drawtdwire{0}{3}{$\headvariableof{\formula}$}
        \drawtensorblock{0}{0}{2}{$\bencodingof{\formula}$}
        \drawtdwire{5}{3}{$\headvariableof{\secformula}$}
        \drawtensorblock{5}{0}{2}{$\bencodingof{\secformula}$}

        \draw (1,-3) to[bend right=-20] (-1,-1);
        \draw (1,-3) to[bend right=20] (4,-1);
        \drawvariabledot{1}{-3}
        \drawtdwire{1}{-3}{$\catvariableof{0}$}
        \draw (4,-3) to[bend right=-20] (1,-1);
        \draw (4,-3) to[bend right=20] (6,-1);
        \drawvariabledot{4}{-3}
        \drawtdwire{4}{-3}{$\catvariableof{\catorder-1}$}
        \drawtextnode{3}{-4.5}{\colorlabelsize $\cdots$}

        \drawtextnode{10}{0}{$=$}

        \drawtensorblock{15}{0}{2}{$\ones$}
        \drawtdwire{13.5}{-1}{$\catvariableof{0}$}
        \drawtextnode{15.5}{-2.5}{\colorlabelsize $\cdots$}
        \drawtdwire{16.5}{-1}{$\catvariableof{\catorder-1}$}
    }
    When knowing for some $\shortcatindicesin$ that $\formulaat{\indexedshortcatvariables}=1$, we conclude that $\secformulaat{\indexedshortcatvariables}=1$.

    \stantikzpic{
        \drawtensorblock{2.5}{4}{3}{$\impbincon$}
        \drawtdwire{0}{3}{$\headvariableof{\formula}$}
        \drawtensorblock{0}{0}{2}{$\bencodingof{\formula}$}
        \drawtdwire{5}{3}{$\headvariableof{\secformula}$}
        \drawtensorblock{5}{0}{2}{$\bencodingof{\secformula}$}

        \draw (1,-3) to[bend right=-20] (-1,-1);
        \draw (1,-3) to[bend right=20] (4,-1);
        \drawvariabledot{1}{-3}
        \drawtdwire{1}{-3}{$\catvariableof{0}$}
        \draw (4,-3) to[bend right=-20] (1,-1);
        \draw (4,-3) to[bend right=20] (6,-1);
        \drawvariabledot{4}{-3}
        \drawtdwire{4}{-3}{$\catvariableof{\catorder-1}$}
        \drawtextnode{3}{-4.5}{\colorlabelsize $\cdots$}

        \drawsmalltensor{1}{-6}{$\onehotmapof{\catindexof{0}}$}
        \drawsmalltensor{4}{-6}{$\onehotmapof{\catindexof{\catorder\shortminus1}}$}

        \drawtextnode{10}{0}{$=$}

        \drawtensorblock{15}{0}{2}{$\ones$}
        \drawtdwire{13.5}{-1}{$\catvariableof{0}$}
        \drawtextnode{15.5}{-2.5}{\colorlabelsize $\cdots$}
        \drawtdwire{16.5}{-1}{$\catvariableof{\catorder-1}$}
        \drawsmalltensor{13.5}{-4}{$\onehotmapof{\catindexof{0}}$}
        \drawsmalltensor{16.5}{-4}{$\onehotmapof{\catindexof{\catorder\shortminus1}}$}


        \drawtextnode{20}{0}{$=\quad 1$}
    }

    We simplify the left hand side using basis calculus to

    \stantikzpic{
        \drawtensorblock{2.5}{4}{3}{$\impbincon$}
        \drawtdwire{0}{3}{$\headvariableof{\formula}$}
        \drawtensorblock{-1}{0}{2.5}{$\onehotmapof{\formulaat{\indexedshortcatvariables}}$}
        \drawtdwire{5}{3}{$\headvariableof{\secformula}$}
        \drawtensorblock{5}{0}{2}{$\bencodingof{\secformula}$}
        \drawtdwire{3.5}{-1}{$\catvariableof{0}$}
        \drawtextnode{5.5}{-2.25}{\colorlabelsize $\cdots$}
        \drawtdwire{6.5}{-1}{$\catvariableof{\catorder\shortminus1}$}
        \drawsmalltensor{3.5}{-4}{$\onehotmapof{\catindexof{0}}$}
        \drawsmalltensor{6.5}{-4}{$\onehotmapof{\catindexof{\catorder\shortminus1}}$}

        \drawtextnode{11}{0}{$= \quad 1$}
    }
    When known that $\formulaat{\indexedshortcatvariables}=1$, this further simplifies to

    \stantikzpic{
        \drawtensorblock{5}{4}{1}{$\tbasis$}
        \drawtdwire{5}{3}{$\headvariableof{\secformula}$}
        \drawtensorblock{5}{0}{2}{$\bencodingof{\secformula}$}
        \drawtdwire{3.5}{-1}{$\catvariableof{0}$}
        \drawtextnode{5.5}{-2.25}{\colorlabelsize $\cdots$}
        \drawtdwire{6.5}{-1}{$\catvariableof{\catorder\shortminus1}$}
        \drawsmalltensor{3.5}{-4}{$\onehotmapof{\catindexof{0}}$}
        \drawsmalltensor{6.5}{-4}{$\onehotmapof{\catindexof{\catorder\shortminus1}}$}

        \drawtextnode{11}{0}{$= \quad 1$}
    }
    which is equal to $\secformulaat{\indexedshortcatvariables}=1$.
\end{example}

More sophisticated reasoning schemes based on contraction equations are applied in \cite{sato_linear_2017,tsilionis_tensor-based_2024}, where recursive linear equations are derived.